
%\needspace{4\baselineskip}
\color{uniblue}{\subsection{Автоматическое повторное включение (АПВ)}}
\color{black}

\begin{enumerate}[label=\arabic{section}.\arabic{subsection}.\arabic*, labelsep=4pt, leftmargin=0pt, itemindent=57pt]

\item
Около 85\% повреждений изоляции воздушных линий связано с дуговыми короткими замыканиями, которые по своей природе имеют неустойчивый характер и исчезают при отключении линии. Это означает, что линия может быть снова включена автоматикой повторного включения. В устройстве автоматическое повторное включение представлено функцией «АПВ».

\item
Логика АПВ реализует следующие возможности:
\begin{itemize}
    \item один цикл АПВ после аварийного отключения выключателя;
    \item формирование блокировок АПВ;
    \item формирование сигнала готовности АПВ;
    \item контроль напряжения шин и присоединения для пуска АПВ;
    \item пуск алгоритма улавливания синхронизма по факту срабатывания АПВ.
\end{itemize}

\item
Готовность АПВ формируется по истечению выдержки времени таймера Т1 при оперативном включении выключателя (наличие активного сигнала <<В: Включен>>), отсутствии сигналов блокировки включения выключателя от логики КСВ, сигналов запрета АПВ и наличие сигнала РФК.

\item
Пусковые условия АПВ формируются при выполнении следующих условий:
\begin{itemize}
    \item аварийное отключение выключателя; 
    \item готовность АПВ; 
    \item возможность проверки синхронизма (программный переключатель SGF3);
    \item контроль напряжения (программный переключатель SGF2).
\end{itemize}

\item
Команда на включение выключателя формируется по истечении выдержки времени таймера Т2 при условии отсутствия блокировки включения выключателя.

\item
Если затянулась команда включения выключателя, то по истечении времени таймера Т3 сформируется выходной сигнал задержки включения выключателя <<АПВ: Задержка включения>> и сбросится готовность АПВ.

\item
Сброс готовности АПВ формируется при наличии одного из следующих условий:
\begin{itemize}
    \item функция АПВ выведена из работы;
    \item отсутствует сигнал РФК или активны сигналы запрета АПВ;
    \item наличие команды <<АПВ: Включить>>;    
    \item сформирован сигнал задержки включения выключателя.
\end{itemize}

\item
Ввод функции в работу на этапе параметрирования устройства осуществляется программным переключателем SGF1 <<Ввод\_функции>>. Информация о введенной функции отображается выходным сигналом <<АПВ: Ввод>>.

\item
Функция может быть выведена из работы оперативно путем активации сигнала <<ОВ АПВ>>, что характеризуется активным состоянием сигнала <<АПВ: Оперативный вывод>> на выходе функции. 
\item
Алгоритм работы функции выполнен в соответствии с рисунком \ref{apv:img1}. В таблице \ref{apv:tbl1} приведены параметры, необходимые для настройки функции.

\vspace{1mm}
\begin{figure}[H]
\centering
\includegraphics[width=1\textwidth,height=1\textheight,keepaspectratio]{\fbpath/305. АПВ ВВ/_latex/img/apv.pdf}
\captionof{figure}{Функциональная схема логики АПВ}\label{apv:img1}
\end{figure}

\end{enumerate}
